% !TeX root = ../despliegue.tex
% ===========================================================================
%  Operación del día a día.
% ===========================================================================

\section{Operación}

\subsection{Arranque limpio}

\begin{comando}
docker compose down -v && docker compose up -d --wait
\end{comando}

El \cmd{-v} borra el volumen \cmd{pgdata}. Es lo que hace que vuelvan a
ejecutarse los scripts de \cmd{infra/postgres/init/} y las migraciones con su
seed, dejando el sistema exactamente como la primera vez.

\aviso{Sin \cmd{-v} los datos se conservan entre arranques: es lo que se quiere
entre sesiones de trabajo, y lo que \textbf{no} se quiere antes de una
demostración, porque las reservas de ejemplo se sembraron con fechas relativas
al primer arranque y con el tiempo quedan en el pasado.}

\subsection{Detener sin borrar}

\begin{comando}
docker compose down     # para los contenedores, conserva los datos
\end{comando}

\subsection{Ver los registros}

\begin{comando}
docker compose logs -f ms-reservas     # uno
docker compose logs -f                 # todos
\end{comando}

\subsection{Reconstruir tras un cambio de código}

\begin{comando}
docker compose up -d --wait --build ms-reservas
\end{comando}

Sin \cmd{--build}, Compose reutiliza la imagen anterior y el cambio no surte
efecto. Es la equivocación más fácil de cometer, y la más desconcertante,
porque el sistema arranca sin error alguno mostrando la versión vieja.

\subsection{Trabajar solo en el frontend}

\begin{comando}
docker compose up shell mf-reservas mf-administracion mf-reportes
\end{comando}

Levanta los cuatro microfrontends sin el backend. Las llamadas a \cmd{/api}
fallarán, que es lo esperado en ese modo.
